%! Author = Omar Iskandarani
%! Date = 4/3/2026
%! Affiliation = Independent Researcher, Groningen, The Netherlands
%! License = © 2025 Omar Iskandarani. All rights reserved. This manuscript is made available for academic reading and citation only. No republication, redistribution, or derivative works are permitted without explicit written permission from the author. Contact: info@omariskandarani.com
%! ORCID = 0009-0006-1686-3961
%! DOI = 10.5281/zenodo.xxx

\documentclass[11pt]{article}
\usepackage[margin=1in]{geometry}
\usepackage[T1]{fontenc}
\usepackage[utf8]{inputenc}
\usepackage{lmodern}
\usepackage{amsmath,amssymb,amsthm,mathtools,bm}
\usepackage{physics}
\usepackage[hidelinks]{hyperref}
\usepackage{microtype}
\usepackage{enumitem}
\usepackage{xcolor}

% --- SST macro prelude ---
\newcommand{\VolH}[1]{\operatorname{Vol}_{\!\mathbb{H}}(#1)}
\newcommand{\swirlarrow}{\mathchoice{\mkern-2mu\scriptstyle\boldsymbol{\circlearrowleft}}{\mkern-2mu\scriptstyle\boldsymbol{\circlearrowleft}}{\mkern-2mu\scriptscriptstyle\boldsymbol{\circlearrowleft}}{\mkern-2mu\scriptscriptstyle\boldsymbol{\circlearrowleft}}}
\newcommand{\swirlarrowcw}{\mathchoice{\mkern-2mu\scriptstyle\boldsymbol{\circlearrowright}}{\mkern-2mu\scriptstyle\boldsymbol{\circlearrowright}}{\mkern-2mu\scriptscriptstyle\boldsymbol{\circlearrowright}}{\mkern-2mu\scriptscriptstyle\boldsymbol{\circlearrowright}}}
\newcommand{\Vol}{\operatorname{Vol}}
\newcommand{\rhof}{\rho_{\!f}}
\newcommand{\rhoE}{\rho_{\!E}}
\newcommand{\rhoM}{\rho_{\!m}}
\newcommand{\rhocore}{\rho_{\text{core}}}
\newcommand{\vswirl}{\mathbf{v}_{\!\boldsymbol{\circlearrowleft}}}
\newcommand{\vswirlcw}{\mathbf{v}_{\!\boldsymbol{\circlearrowright}}}
\newcommand{\SwirlClock}{S_t^{\boldsymbol{\circlearrowleft}}}
\newcommand{\SwirlClockcw}{S_t^{\boldsymbol{\circlearrowright}}}
\newcommand{\vnorm}{\lVert \vswirl \rVert}
\providecommand{\rc}{r_c}

\title{SST-74: Defect-Bound Chirality Probes and Near-Zero Modes\\
\large A Translation Note from Impurity Spectroscopy in a Chiral Kagome Superconductor}
\author{Omar Iskandarani}
\date{\today}

\begin{document}
    \maketitle

    \begin{abstract}
        We formulate a narrow Swirl-String Theory (SST) translation of a recent impurity-spectroscopy result in the kagome superconductor CsV$_3$Sb$_5$. The source paper reports two central features: first, a defect-localized resonance with a strongly mirror-broken spatial pattern interpreted as evidence for hidden chiral order; second, a robust zero-bias peak near a vacancy, discussed as a defect-bound near-zero mode associated with a sign-changing order-parameter boundary. The present note does not import the condensed-matter ontology literally. Instead, it extracts the transferable structural content and rewrites it in SST terms. The main translatable claim is that a localized defect can reveal hidden handed order of the host medium through a symmetry-broken localized response, and that a sign-changing branch field near a defect can support a near-zero localized mode. We state explicitly which parts translate cleanly, which remain analogical only, and which are currently unknown in SST. A minimal effective SST operator is proposed for numerical testing.
    \end{abstract}

    \section{Purpose and scope}
        This note is intentionally narrow. It is not a claim that Kondo physics, Yu--Shiba--Rusinov (YSR) states, Bogoliubov quasiparticles, or Majorana zero modes are already present in canonical SST. The goal is more modest:
        \begin{enumerate}[label=\arabic*.]
    \item identify the structurally transferable content of the source paper,
    \item rewrite that content in SST-compatible language,
    \item state explicitly where the translation stops, and
    \item propose a minimal SST defect-mode ansatz suitable for falsifiable numerical work.
        \end{enumerate}

        The translation criterion used throughout is strict:
        \begin{quote}
            A paper element is called \emph{transferred} only if it can be rewritten as an SST-level field, defect, operator, or observable without silently importing orthodox condensed-matter ontology.
        \end{quote}

    \section{Source-paper kernel}
        The source paper studies impurity-induced local electronic states in the kagome superconductor CsV$_3$Sb$_5$. Its two central observational claims may be summarized as follows:
        \begin{enumerate}[label=(\alph*)]
    \item Cr dopants produce a localized resonance whose spatial pattern breaks all in-plane mirror symmetries, interpreted as evidence for hidden chiral order in the host phase.
    \item V vacancies produce a robust zero-bias peak inside the superconducting gap, interpreted as a defect-bound near-zero mode consistent with a sign-changing order-parameter boundary.
        \end{enumerate}

        For SST, the transferable core is therefore not the specific condensed-matter vocabulary but the following narrower mechanism:
        \begin{equation}
            \boxed{
                \text{A localized defect can reveal hidden handed order of the host medium through a symmetry-broken localized response.}
            }
        \end{equation}
        A second transferable kernel is:
        \begin{equation}
            \boxed{
                \text{A sign-changing branch field near a defect can support a localized near-zero mode.}
            }
        \end{equation}

    \section{SST translation dictionary}
        We separate the translation into three categories.

        \subsection{Transferred}
            \begin{enumerate}[label=\arabic*.]
    \item \textbf{Defect as probe of hidden order.} A localized defect in a handed background can act as a spectroscopic probe of the background handedness.
    \item \textbf{Mirror-broken local response as chirality diagnostic.} If the defect couples to a pseudoscalar or parity-odd structure of the host medium, the localized response need not preserve mirror symmetry.
    \item \textbf{Near-zero defect mode from branch reversal.} A sign-changing or branch-changing order field around a defect can produce a localized low-frequency mode.
            \end{enumerate}

        \subsection{Analogical only}
            \begin{enumerate}[label=\arabic*.]
    \item \textbf{Kondo screening.} There is currently no canonical SST object corresponding to a Fermi-sea screening cloud around a magnetic impurity.
    \item \textbf{Majorana mode language.} A strict Bogoliubov--de Gennes particle-hole Majorana interpretation is not currently part of SST.
    \item \textbf{Topological superconducting surface state.} This phrase remains specific to band topology and superconducting order, not directly to canonical SST.
            \end{enumerate}

        \subsection{Unknown / not yet translatable}
            \begin{enumerate}[label=\arabic*.]
    \item Which precise SST field should replace the sign-changing superconducting order parameter.
    \item Which microscopic SST operator should replace the defect--chirality coupling of the source paper.
    \item Whether the appropriate SST near-zero mode should be interpreted as torsional, swirl-clock, or another branch excitation.
            \end{enumerate}

    \section{Minimal SST fields}
        To express the translated mechanism, introduce the following coarse-grained fields.

        First, a handedness field
        \begin{equation}
            \chi(\mathbf{x}) \in \{-1,+1\},
        \end{equation}
        which distinguishes two local chirality sectors of the medium.

        Second, a branch or sign field
        \begin{equation}
            \sigma(\mathbf{x}) \in \{-1,+1\},
        \end{equation}
        which represents a local branch assignment of an order field that may reverse across a defect-centered boundary.

        Third, a localized defect potential
        \begin{equation}
            V_{\mathrm{def}}(\mathbf{x}) = V_0\,f(\mathbf{x}-\mathbf{x}_0),
        \end{equation}
        with center $\mathbf{x}_0$ and spatial scale $\ell$.

        A natural SST coarse-grained pseudoscalar associated with local handedness is the helicity-density-like quantity
        \begin{equation}
            h(\mathbf{x}) = \mathbf{u}(\mathbf{x})\cdot\bigl(\nabla\times \mathbf{u}(\mathbf{x})\bigr),
        \end{equation}
        where $\mathbf{u}$ is an effective medium velocity field. At phenomenological level one may let
        \begin{equation}
            \chi(\mathbf{x}) = \operatorname{sgn}\!\bigl(h(\mathbf{x})\bigr),
        \end{equation}
        although this identification is not yet canonical.

    \section{Minimal effective operator}
        A minimal translated SST operator that captures both chirality-sensitive mirror breaking and defect-bound near-zero modes is
        \begin{equation}
            \left[
                -\partial_t^2
            + c_T^2 \nabla^2
            - m_0^2
            - m_1^2 \, \sigma(\mathbf{x})
                + \lambda_\chi \, \chi(\mathbf{x})\, \hat{\mathbf{n}}\cdot(\nabla\times)
                + V_{\mathrm{def}}(\mathbf{x})
            \right]\psi = 0.
            \label{eq:min-op}
        \end{equation}
        Here:
        \begin{itemize}
            \item $c_T$ is an effective propagation speed for the relevant localized mode sector,
            \item $m_0$ is a baseline mode gap,
            \item $m_1^2\sigma(\mathbf{x})$ is a branch-sensitive local mass term,
            \item $\lambda_\chi$ measures the strength of chirality-sensitive coupling,
            \item $\hat{\mathbf{n}}$ is a preferred local orientation set by the background structure,
            \item $V_{\mathrm{def}}$ is the defect-localizing perturbation.
        \end{itemize}

        Equation \eqref{eq:min-op} is \emph{not} derived from the source paper. It is a translation ansatz that captures the transferable operator-level structure.

    \section{Mirror-broken localized response}
        Suppose the background has nonzero handedness,
        \begin{equation}
            \chi(\mathbf{x}) \neq 0,
        \end{equation}
        and that the defect couples to the parity-odd term
        \begin{equation}
            \lambda_\chi\,\chi(\mathbf{x})\,\hat{\mathbf{n}}\cdot(\nabla\times).
        \end{equation}
        Then the local mode intensity
        \begin{equation}
            I(\mathbf{x};\mathbf{x}_0)=\abs{\psi(\mathbf{x};\mathbf{x}_0)}^2
        \end{equation}
        need not respect the mirror symmetries that would otherwise be expected from the surrounding geometry. In other words, localized defect spectroscopy can become a chirality diagnostic:
        \begin{equation}
            \chi \neq 0 \,\wedge\, \lambda_\chi\neq 0
            \quad\Longrightarrow\quad
            I(\mathbf{x}) \text{ may be mirror-broken.}
        \end{equation}

        This is the clean SST translation of the source paper's first central motif.

    \section{Near-zero mode from branch reversal}
        The second motif is the appearance of a robust near-zero localized mode when an order field changes sign near a defect. The structurally standard one-dimensional prototype is
        \begin{equation}
            \left[-\partial_x^2 + m^2(x)\right]\psi(x) = \omega^2\psi(x),
            \label{eq:1d-bound}
        \end{equation}
        with an odd mass profile
        \begin{equation}
            m(-x)=-m(x).
        \end{equation}
        If the effective local gap reverses sign across a defect-centered interface, Equation \eqref{eq:1d-bound} can support a mode with small $\omega$ localized near the interface.

        In translated SST language, the sign-changing field is represented by $\sigma(\mathbf{x})$. The defect-centered radial prototype is
        \begin{equation}
            \sigma(r)=
            \begin{cases}
                -1, & r<r_0,\\
                +1, & r>r_0.
            \end{cases}
        \end{equation}
        Substituting this into Equation \eqref{eq:min-op} gives a direct SST test model for a defect-bound near-zero mode.

    \section{What remains unknown}
        The translation remains incomplete in three important ways.

        \subsection{Unknown replacement for the superconducting sign field}
            The source paper uses a sign-changing superconducting order parameter. In SST, the correct counterpart is not yet fixed. Plausible candidates are:
            \begin{equation}
                Q(\mathbf{x}) = \chi(\mathbf{x}),
                \qquad
                Q(\mathbf{x}) = \Theta(\mathbf{x})-\Theta_0,
                \qquad
                Q(\mathbf{x}) = \SwirlClock(\mathbf{x})-\SwirlClock_0,
            \end{equation}
            but none of these identifications is presently canonical.

        \subsection{Unknown microscopic coupling}
            The paper motivates a coupling between defect-localized response and hidden chirality. SST does not yet possess a derived microscopic operator for that coupling. The term
            \begin{equation}
                \lambda_\chi \, \chi(\mathbf{x})\, \hat{\mathbf{n}}\cdot(\nabla\times)
            \end{equation}
            is therefore only an effective phenomenological proxy.

        \subsection{Unknown mode identity}
            Even if a near-zero bound state is found in Equation \eqref{eq:min-op}, its interpretation is not yet determined. It could be:
            \begin{enumerate}[label=(\alph*)]
    \item a torsional bound mode,
    \item a swirl-clock branch mode,
    \item a mixed chirality--branch excitation,
    \item or a more complicated defect-pinned collective mode.
            \end{enumerate}

    \section{Falsifiable numerical program}
        A direct numerical test can be built around Equation \eqref{eq:min-op} in two spatial dimensions. Let
        \begin{equation}
            V_{\mathrm{def}}(r)=V_0\exp\!\left(-\frac{r^2}{\ell^2}\right),
        \end{equation}
        and define trial profiles
        \begin{equation}
            \chi(x,y)=\chi_0\tanh\!\left(\frac{x}{w_\chi}\right),
            \qquad
            \sigma(x,y)=\tanh\!\left(\frac{x}{w_\sigma}\right).
        \end{equation}
        Then the following predictions are sharp:
        \begin{align}
            \chi_0=0 &\Longrightarrow \text{the defect pattern should remain more mirror-symmetric},\\
            \chi_0\neq 0 &\Longrightarrow \text{the defect pattern may become mirror-broken},\\
            \sigma \text{ sign reversal} &\Longrightarrow \text{a localized low-frequency mode becomes more likely.}
        \end{align}

        A practical observable is the local density
        \begin{equation}
            I(x,y)=\abs{\psi(x,y)}^2,
        \end{equation}
        compared across parameter scans in $(\lambda_\chi, m_1, V_0, \ell, w_\chi, w_\sigma)$.

    \section{Relation to current SST ontology}
        This note fits naturally with the current SST direction in which bosons are rings or simple strings and the photon is treated as a torsional pulse rather than a static knot. In that setting, a defect-bound near-zero mode is plausibly more naturally interpreted as a localized torsional or branch excitation than as a matter-knot excitation. That interpretation remains tentative, but it is ontologically closer to the present SST picture than a literal import of condensed-matter quasiparticles.

    \section{Main conclusion}
        The source paper does not provide direct SST mathematics in the strict sense. It does, however, provide a strong and useful structural lesson:
        \begin{equation}
            \boxed{
                \text{defects can diagnose hidden handed order, and sign-changing branch structure can trap near-zero modes.}
            }
        \end{equation}
        That lesson translates well into SST and motivates a concrete effective operator, Equation \eqref{eq:min-op}, for numerical exploration. The translation is therefore partial but substantial:
        \begin{equation}
            \boxed{\text{philosophical transfer: strong}},
            \qquad
            \boxed{\text{mathematical transfer: moderate and structural}},
            \qquad
            \boxed{\text{literal ontology transfer: incomplete}}.
        \end{equation}

        \paragraph{Child-level analogy.}
            A tiny bump on a spinning trampoline can reveal which hidden way the whole fabric prefers to twist, and if the twist changes sign across the bump, a small trapped wiggle can live there.

            \begin{thebibliography}{99}

                \bibitem{Tu2025arXiv}
                Y.~Tu, Z.~Zhang, W.~Lu, T.~Han, R.~Lv, Z.~Wang, Z.~Zhou, X.~Hou, N.~Hao, Z.~Wang, X.~Chen, and L.~Shan,
                \newblock Symmetry-Broken Kondo Screening and Zero-Energy Mode in the Kagome Superconductor CsV$_3$Sb$_5$,
                \newblock arXiv preprint arXiv:2502.20733, 2025.
                \newblock DOI: \url{https://doi.org/10.48550/arXiv.2502.20733}.

                \bibitem{Tu2026NatPhys}
                Y.~Tu, Z.~Zhang, W.~Lu, Y.~Xiao, T.~Han, R.~Lv, Z.~Wang, Z.~Zhou, X.~Hou, N.~Hao, X.~Chen, and L.~Shan,
                \newblock Symmetry-broken Kondo screening and zero-energy mode in a kagome superconductor,
                \newblock \emph{Nature Physics}, published online 24 March 2026.
                \newblock DOI: \url{https://doi.org/10.1038/s41567-026-03223-5}.

                \bibitem{Christensen2022LoopCurrents}
                M.~H.~Christensen, T.~Birol, B.~M.~Andersen, and R.~M.~Fernandes,
                \newblock Loop currents in $A$V$_3$Sb$_5$ kagome metals: Multipolar and toroidal magnetic orders,
                \newblock \emph{Physical Review B} \textbf{106}, 144504 (2022).
                \newblock DOI: \url{https://doi.org/10.1103/PhysRevB.106.144504}.

                \bibitem{Holbaek2023Disorder}
                S.~C.~Holb{\ae}k, M.~H.~Christensen, A.~Kreisel, and B.~M.~Andersen,
                \newblock Unconventional superconductivity protected from disorder on the kagome lattice,
                \newblock \emph{Physical Review B} \textbf{108}, 144508 (2023).
                \newblock DOI: \url{https://doi.org/10.1103/PhysRevB.108.144508}.

                \bibitem{Madhavan1998KondoSTM}
                V.~Madhavan, W.~Chen, M.~F.~Crommie, N.~S.~Wingreen, and H.~C.~Manoharan,
                \newblock Tunneling into a Single Magnetic Atom: Spectroscopic Evidence of the Kondo Resonance,
                \newblock \emph{Science} \textbf{280}, 567--569 (1998).
                \newblock DOI: \url{https://doi.org/10.1126/science.280.5363.567}.

            \end{thebibliography}

\end{document}